Hungarian astronomers significantly contributed to the study of pulsating
variable stars of RR Lyrae type, which show cyclic amplitude modulation in
their light variation (Blazhko effect).

The light variation of an RR Lyrae star is observed at two different
wavelengths: $\lambda_1 = 500$\,nm and $\lambda_2 = 2000$\,nm. At each
wavelength, we see into the star at a different depth. We refer to the depths
as layer 1 and layer 2. One can approximate the light intensity of a star at
a given wavelength by the radius and the black-body radiation intensity at
the appropriate depth. Moreover, the Wien approximation can be used for the
black-body radiation to calculate the emitted power from unit surface area:
\[ F(\lambda, T) \propto \frac{1}{\lambda^5}
   \exp\left(-\frac{hc}{k\lambda T}\right), \]
where $h$ is the Planck constant, $k$ is the Boltzmann constant, and $c$ is
the speed of light. To simplify the calculations we introduce a new constant
$C_b = hc/k \approx 0.0144$\,m\,K.

\begin{description}
\item[(a)] Assume that the temperature varies between $T_1 = 6000$\,K and
  $T_2 = 7400$\,K in each of the layers and ignoring radius variation, what
  is the ratio of the amplitudes of variation in magnitudes at the two
  wavelengths? \hfill[5\,p]
\item[(b)] What is the peak to peak amplitude of the light curve at
  $\lambda_1$? Use the magnitude scale. \hfill[3\,p]
\item[(c)] Ignoring temperature variation, what is the contribution of the
  radius variation to the light-curve amplitude for a given wavelength, if
  $R_{\min} = 0.9\,\langle R\rangle$ and
  $R_{\max} = 1.05\,\langle R\rangle$? $\langle R\rangle$ is the mean radius
  of the given layer. \hfill[3\,p]
\item[(d)] Recent observations and models show that the radius of the
  photosphere is only minimally modulated during the Blazhko cycle; however,
  the temperature variation is significant. As a result, the amplitude of
  light curve itself keeps varying. Let us assume that during the minimum
  amplitude of pulsation the temperature variation is reduced to
  $T_{\min} = 6100$\,K and $T_{\max} = 6900$\,K. What is the modulation
  amplitude at the two wavelengths? For the maximum amplitude, use the
  temperature values given in part (a). \hfill[5\,p]
\item[(e)] Which statement is correct? (There can be more than one correct
  selection.) Mark the letter of your answer with $\times$ on the answer
  sheet. \hfill[4\,p]
  \begin{description}
  \item[(A)] It is easier to observe the Blazhko effect in the infrared
    band.
  \item[(B)] The temperature variation dominates the visual light curve.
  \item[(C)] If we neglect the radius variation, the amplitude is inversely
    proportional to the wavelength.
  \item[(D)] Multi-colour observations are not useful in understanding the
    Blazhko effect.
  \end{description}
\end{description}
